\section{Discussion}
\label{sec:discussion}

This paper is one application of the masked-cause coarsening
framework of \citet{towell2026masked}; the shared identifiability
vocabulary is what the series argues for, and each application
contributes a domain-specific construction (here, the glass-ceiling
symmetry and the gold-set sample-complexity bound).

\subsection{Relation to prior work on label-model identifiability}

Label-model identifiability has a deep prior literature, and an
honest account places this paper's contribution against it.

\paragraph{Dawid--Skene is the ancestor.} \citet{dawid1979maximum}
posed and solved the core problem: estimate per-rater error rates and
a class prior, by maximum likelihood via EM, from rater votes alone.
Every label model in programmatic weak supervision, and the
masked-cause statement of it in this paper, is a descendant of that
work. The EM solver of Dawid--Skene is the maximum-likelihood solver
for the coarsened likelihood \eqref{eq:bg-c1c2c3}; the framework here
does not replace it, it gives it an identifiability vocabulary.

\paragraph{Moment-based identifiability is established.} Data
programming \citep{ratner2016data} showed that LF accuracies are
recoverable from agreement statistics under conditional independence;
the triplet method of \citet{fu2020fast} made the recovery a
closed form in third-order agreement moments; the multi-task
extension \citep{ratner2019training,ratner2020training} handled
structured dependence with a known dependency graph;
\citet{bach2017learning} showed the dependency structure can
sometimes be learned, but only under restrictive conditions.
\Cref{thm:ci-identifiability} is a re-derivation of this body of
results, not a new identifiability claim. The crowdsourcing
literature \citep{karger2011iterative,zhang2016spectral} solved the
same estimation problem for human annotators; the spectral and
tensor-decomposition tools \citep{anandkumar2014tensor} use the
same low-rank moment structure that underlies the masked-cause
augmented-matrix rank condition (\cref{thm:bg-id}). The Snorkel
ecosystem \citep{ratner2018snorkelmetal,bach2019snorkeldrybell,
sala2019multiresolution,zhang2022survey} has expanded weak
supervision to multi-task, deployment, sequential, and interactive
\citep{boecking2021interactive} settings; the masked-cause framework
applies to each, because all share the LF-vote-matrix input.

\paragraph{What this paper adds.} The contribution is narrower than
``label-model identifiability'' and is fourfold. (i) The masked-cause
unification: LF votes are a candidate-set coarsening, abstention is a
non-informative candidate set, a confident vote is a narrowed one,
and gold labels are singletons. This places weak supervision in the
same frame as the other applications of the masked-cause framework
\citep{towell2026masked}. (ii) The C1--C2--C3 classification of LF
ensembles: the three coarsening-at-random conditions become checkable
statements about an LF set (support, label-only error dependence,
parameter separation). (iii) The explicit glass-ceiling construction:
the accuracy-complement symmetry of \cref{thm:glass-ceiling} is a
precise statement, with an exact label-flip witness, of why gold-free
estimation cannot orient the model. (iv) The gold-set
sample-complexity theorem for the dependent case
(\cref{thm:goldset}): prior work establishes the conditionally
independent corner ($r = 0$); the $\Theta(r / \mathrm{gap}^2)$ bound
for total model recovery extends it to the correlated ensembles that
practice produces, and the $1/\mathrm{gap}^2$ scaling is confirmed in
simulation. We do not
claim to invent label-model identifiability; we claim a unifying
frame, a classification, and a sample-complexity result for the
dependent case.

\subsection{What the framework does \emph{not} address}

\begin{itemize}
\item \textbf{Multiclass and structured outputs.} The development is
  binary. The masked-cause framework extends to a categorical latent
  cause directly, and Dawid--Skene is natively multiclass; the
  C1--C2--C3 statements and the rank arguments carry over, but we did
  not write out the multiclass theorems here.

\item \textbf{Unknown dependency structure.} \Cref{thm:goldset}
  assumes the rank deficit $r$ is a property of the ensemble, but
  $r$ is itself estimated. Recovering the dependency graph from data,
  rather than assuming it as in \citet{ratner2019training}, is a
  harder problem the framework does not solve.

\item \textbf{Abstention that depends on the label.} We model
  coverage $\beta_j$ as independent of $Y$. An LF that abstains
  preferentially on one class is a label-dependent coarsening, a C2
  violation the present analysis does not cover; the C2-relaxation
  results of \citet{towell2026mdrelax} are the natural tool.

\item \textbf{The downstream model.} The framework analyzes the label
  model. Whether the probabilistic training labels it emits yield a
  good downstream classifier is a separate generalization question,
  governed by the end-to-end analysis of \citet{ratner2016data}.

\item \textbf{Adversarial or systematically wrong LFs.} C1 fails when
  an LF is confidently wrong beyond the modeled noise. Such LFs need
  detection and removal, which the framework flags as a C1 violation
  but does not itself perform.
\end{itemize}

\subsection{Broader implications}

The framework shifts the question a weak-supervision practitioner
asks from ``which label model should I use?'' to ``which
coarsening-at-random conditions does my LF ensemble satisfy?'' If the
LFs are conditionally independent and better than random, the
ensemble is on the identifiable side of the boundary and a gold-free
pipeline is sound (\cref{thm:ci-identifiability}). If the LFs are
correlated, the ensemble is rank-deficient and gold labels are
required, with a budget set by the rank deficit and the accuracy
margin (\cref{thm:goldset}). If an LF is confidently wrong, C1 fails
and no aggregation rescues it. This is the same assumption-specific
reasoning previously brought to other masked-cause applications in
the framework series \citep{towell2026masked}: recognizing the
common masked-cause structure lets one identifiability apparatus
serve measurement problems that look unrelated on the surface. The
framework series collectively argues that the candidate-set
abstraction is a general tool for any inference problem in which
observations partially-determine a latent cause.
